\documentclass{article}
\usepackage{amsmath,amssymb,amsthm}
\usepackage{algorithm,algorithmic}
\usepackage{bm}
\usepackage{hyperref}
\newtheorem{theorem}{Theorem}
\newtheorem{lemma}{Lemma}
\newtheorem{assumption}{Assumption}
\newcommand{\E}{\mathbb{E}}
\newcommand{\R}{\mathbb{R}}

\begin{document}

\section*{Algorithm Name}
\textbf{Randomized Block Coordinate Descent (RBCD)}  

At each iteration a single block of coordinates is sampled uniformly at random and updated by a gradient step while all other blocks remain unchanged.

\vspace{0.5em}
\section*{Mathematical Formulation}
Let the decision variable be partitioned into $B$ blocks
\[
\bm w = \big(w^{(1)},w^{(2)},\dots ,w^{(B)}\big),\qquad 
w^{(i)}\in\R^{d_i},\;\;\sum_{i=1}^{B}d_i = d .
\]

Given a convex differentiable objective $f:\R^{d}\to\R$, the algorithm proceeds as follows.
\begin{enumerate}
    \item Initialise $\bm w^{0}\in\R^{d}$.
    \item For $t=0,1,2,\dots$:
    \begin{enumerate}
        \item Sample a block index $i_t\in\{1,\dots ,B\}$ uniformly:
        \[
        \mathbb{P}(i_t=i)=\frac{1}{B},\qquad i=1,\dots ,B .
        \]
        \item Compute the partial gradient w.r.t. block $i_t$:
        \[
        \bm g^{(i_t)}_t \;:=\; \nabla_{i_t} f\!\big(\bm w^{t}\big)
        \in \R^{d_{i_t}} .
        \]
        \item Update only the selected block with a stepsize $\eta>0$:
        \[
        w^{(i_t)}_{t+1}= w^{(i_t)}_{t} - \eta\,\bm g^{(i_t)}_t,
        \qquad
        w^{(j)}_{t+1}= w^{(j)}_{t}\;\;\forall j\neq i_t .
        \]
    \end{enumerate}
\end{enumerate}
The algorithm stops after a pre‑specified number of iterations $T$ or when a
convergence criterion is satisfied.

\vspace{0.5em}
\section*{Pseudocode}
\begin{algorithm}[H]
\caption{Randomized Block Coordinate Descent}
\begin{algorithmic}[1]
\STATE \textbf{Input:} stepsize $\eta>0$, number of iterations $T$, initial point $\bm w^{0}$.
\FOR{$t=0$ \TO $T-1$}
    \STATE Sample $i_t\sim \text{Uniform}\{1,\dots ,B\}$.
    \STATE Compute partial gradient $\bm g^{(i_t)}_t \gets \nabla_{i_t}f(\bm w^{t})$.
    \STATE Update selected block:
    \[
    w^{(i_t)}_{t+1}\gets w^{(i_t)}_{t}-\eta\,\bm g^{(i_t)}_t .
    \]
    \STATE For all $j\neq i_t$, set $w^{(j)}_{t+1}\gets w^{(j)}_{t}$.
\ENDFOR
\STATE \textbf{Output:} $\bm w^{T}$.
\end{algorithmic}
\end{algorithm}

\vspace{0.5em}
\section*{Dry Run Example}
Consider the one‑dimensional problem
\[
f(w) = (w-3)^2 ,\qquad w\in\R .
\]
Here $B=1$, $L_1=2$ (since $f$ is $2$‑smooth) and the full gradient coincides with the block gradient.
Take $\eta=0.1$, $w^{0}=0$ and run three iterations.

\begin{center}
\begin{tabular}{c|c|c|c}
$t$ & $w^{t}$ & $\nabla f(w^{t})$ & $f(w^{t})$ \\ \hline
0 & $0$ & $2(0-3) = -6$ & $9$\\
1 & $0-0.1(-6)=0.6$ & $2(0.6-3)=-4.8$ & $(0.6-3)^2=5.76$\\
2 & $0.6-0.1(-4.8)=1.08$ & $2(1.08-3)=-3.84$ & $(1.08-3)^2=3.6864$\\
3 & $1.08-0.1(-3.84)=1.464$ & $2(1.464-3)=-3.072$ & $(1.464-3)^2=2.366\\
\end{tabular}
\end{center}
The objective value decreases monotonically, illustrating the effect of the random (here deterministic because $B=1$) block update.

\vspace{1em}
\section*{Assumptions}
\begin{assumption}[Convexity]\label{as:convex}
$f:\R^{d}\to\R$ is convex, i.e.
\[
f(\theta \bm x+(1-\theta)\bm y)\le \theta f(\bm x)+(1-\theta)f(\bm y),\qquad 
\forall \bm x,\bm y\in\R^{d},\;\theta\in[0,1].
\]
\end{assumption}
\begin{assumption}[Block‑coordinate $L_i$‑smoothness]\label{as:smooth}
For each block $i\in\{1,\dots ,B\}$ there exists $L_i>0$ such that for any 
$\bm x,\bm y$ that differ only in block $i$,
\[
f(\bm y)\le f(\bm x)+\langle \nabla_i f(\bm x),\; \bm y^{(i)}-\bm x^{(i)}\rangle
+\frac{L_i}{2}\|\bm y^{(i)}-\bm x^{(i)}\|^{2}.
\tag{1}
\]
\end{assumption}
Define the maximum block smoothness constant
\[
L_{\max}\;:=\;\max_{i\in\{1,\dots ,B\}} L_i .
\tag{2}
\]

\begin{assumption}[Existence of a minimiser]\label{as:opt}
There exists at least one optimal point $\bm w^{\star}\in\arg\min_{\bm w}f(\bm w)$.
\end{assumption}

All constants $\eta$, $L_i$ are assumed known; the stepsize will be chosen as
$\eta = 1/L_{\max}$ unless otherwise stated.

\vspace{0.5em}
\section*{Main Convergence Theorem}
\begin{theorem}[Expected suboptimality of RBCD]\label{thm:main}
Let Assumptions~\ref{as:convex}–\ref{as:opt} hold and choose a constant stepsize
\[
\eta = \frac{1}{L_{\max}} .
\tag{3}
\]
Denote by $\{\bm w^{t}\}_{t\ge 0}$ the iterates generated by the algorithm and let
$T\ge 1$.
Then the expected function value satisfies
\[
\E\!\big[\,f(\bm w^{T})-f(\bm w^{\star})\,\big]
\;\le\;
\frac{B\,L_{\max}\,\|\bm w^{0}-\bm w^{\star}\|^{2}}{2\,T}.
\tag{4}
\]
Consequently, RBCD attains an $O(1/T)$ convergence rate in expectation for
convex smooth objectives.
\end{theorem}

\vspace{0.5em}
\section*{Theoretical Properties}
\begin{itemize}
\item \textbf{Stability.} The stepsize $\eta=1/L_{\max}$ guarantees that each
individual block update never increases the objective (see Lemma~\ref{lem:descent}).
\item \textbf{Hyperparameter sensitivity.} The bound (4) depends linearly on
the number of blocks $B$ and on $L_{\max}$. Using a smaller stepsize
$\eta<1/L_{\max}$ yields a larger constant but still preserves the $O(1/T)$
rate.
\item \textbf{Comparison to full‑gradient GD.} Full gradient descent with stepsize
$1/L_{\max}$ yields the identical bound (4) with $B=1$. Hence randomised block
selection introduces at most a factor $B$ in the constant, while each
iteration is cheaper by a factor of roughly $B$.
\end{itemize}

\vspace{0.5em}
\section*{Discussion}
The theorem shows that even though only a single block is updated per iteration,
the expected decrease in objective value matches the classical $1/T$ rate of
full gradient methods up to a multiplicative factor $B$. When $f$ is additionally
strongly convex, the linear rate $O\big((1-\mu/(B L_{\max}))^{T}\big)$ can be
derived (omitted for brevity).

\vspace{1em}
\section*{Preliminaries (Definitions and Lemmas)}
\begin{lemma}[Block descent lemma]\label{lem:descent}
Under Assumption~\ref{as:smooth} and for any stepsize $\eta\le 1/L_i$,
\[
f\!\big(\bm w^{t+1}\big)
\;\le\;
f\!\big(\bm w^{t}\big)
-\frac{\eta}{2}\,\|\nabla_{i_t} f(\bm w^{t})\|^{2}.
\tag{5}
\]
\end{lemma}
\begin{proof}
Condition (1) applied with $\bm x=\bm w^{t}$ and $\bm y$ equal to $\bm w^{t+1}$
(which differs only in block $i_t$) gives
\[
f(\bm w^{t+1})\le f(\bm w^{t})
+\big\langle\nabla_{i_t}f(\bm w^{t}),\,
w^{(i_t)}_{t+1}-w^{(i_t)}_{t}\big\rangle
+\frac{L_{i_t}}{2}\|w^{(i_t)}_{t+1}-w^{(i_t)}_{t}\|^{2}.
\]
Using the update $w^{(i_t)}_{t+1}=w^{(i_t)}_{t}-\eta\nabla_{i_t}f(\bm w^{t})$,
the inner product term becomes $-\eta\|\nabla_{i_t}f(\bm w^{t})\|^{2}$ and the
quadratic term becomes $\frac{L_{i_t}}{2}\eta^{2}\|\nabla_{i_t}f(\bm w^{t})\|^{2}$.
Hence
\[
f(\bm w^{t+1})\le f(\bm w^{t})
-\eta\|\nabla_{i_t}f(\bm w^{t})\|^{2}
+\frac{L_{i_t}\eta^{2}}{2}\|\nabla_{i_t}f(\bm w^{t})\|^{2}.
\]
If $\eta\le 1/L_{i_t}$ then $L_{i_t}\eta^{2}\le\eta$, and thus
\[
f(\bm w^{t+1})\le f(\bm w^{t})
-\frac{\eta}{2}\|\nabla_{i_t}f(\bm w^{t})\|^{2},
\]
which is exactly (5). ∎
\end{proof}

\begin{lemma}[Expected squared block gradient]\label{lem:expgrad}
For any iterate $\bm w^{t}$,
\[
\E_{i_t}\big[\|\nabla_{i_t} f(\bm w^{t})\|^{2}\big]
\;=\;
\frac{1}{B}\sum_{i=1}^{B}\|\nabla_{i} f(\bm w^{t})\|^{2}.
\tag{6}
\]
\end{lemma}
\begin{proof}
Since $i_t$ is sampled uniformly,
\[
\E_{i_t}\big[\|\nabla_{i_t} f(\bm w^{t})\|^{2}\big]
=\sum_{i=1}^{B}\frac{1}{B}\,\|\nabla_{i} f(\bm w^{t})\|^{2}
=\frac{1}{B}\sum_{i=1}^{B}\|\nabla_{i} f(\bm w^{t})\|^{2}.
\tag{7}
\]
∎
\end{proof}

\begin{lemma}[Relation between full gradient norm and function suboptimality]\label{lem:gradbound}
If $f$ is convex and block‑smooth with constants $L_i$, then for any $\bm w$
\[
\frac{1}{2L_{\max}}\|\nabla f(\bm w)\|^{2}
\;\le\;
f(\bm w)-f(\bm w^{\star}).
\tag{8}
\]
\end{lemma}
\begin{proof}
Convexity implies for any $\bm w$ and optimal $\bm w^{\star}$,
\[
f(\bm w)-f(\bm w^{\star})\ge \langle \nabla f(\bm w),\,
\bm w-\bm w^{\star}\rangle .
\tag{9}
\]
By Cauchy–Schwarz,
\[
\langle \nabla f(\bm w),\,
\bm w-\bm w^{\star}\rangle
\ge \frac{1}{\|\bm w-\bm w^{\star}\|}\,
\|\nabla f(\bm w)\|^{2}.
\tag{10}
\]
Smoothness (1) applied with $\bm y=\bm w-\frac{1}{L_{\max}}\nabla f(\bm w)$
and using that each block satisfies the same $L_{\max}$ gives
\[
f\!\Big(\bm w-\frac{1}{L_{\max}}\nabla f(\bm w)\Big)
\le f(\bm w)-\frac{1}{2L_{\max}}\|\nabla f(\bm w)\|^{2}.
\tag{11}
\]
Since $f(\bm w^{\star})\le f\big(\bm w-\frac{1}{L_{\max}}\nabla f(\bm w)\big)$,
combine (11) with $f(\bm w)-f(\bm w^{\star})\ge\frac{1}{2L_{\max}}\|\nabla f(\bm w)\|^{2}$,
which yields (8). ∎
\end{proof}

\vspace{0.5em}
\section*{Main Proof (Step‑by‑Step Derivation)}
We prove Theorem~\ref{thm:main}. Throughout the proof, expectations are taken
with respect to the randomness of the block indices $\{i_0,\dots,i_{t-1}\}$.

\noindent\textbf{Step 1.}  Start from the descent inequality (5) of Lemma~\ref{lem:descent}
with the stepsize (3).  
\[
f(\bm w^{t+1})
\le f(\bm w^{t})-\frac{\eta}{2}\|\nabla_{i_t}f(\bm w^{t})\|^{2}.
\tag{12}
\]

\noindent\textbf{Step 2.}  Take expectation over the random choice of $i_t$ conditional on $\bm w^{t}$.  
Using Lemma~\ref{lem:expgrad} we obtain
\[
\E_{i_t}\!\big[f(\bm w^{t+1})\mid \bm w^{t}\big]
\le f(\bm w^{t})-\frac{\eta}{2}\,
\frac{1}{B}\sum_{i=1}^{B}\|\nabla_{i}f(\bm w^{t})\|^{2}.
\tag{13}
\]

\noindent\textbf{Step 3.}  Relate the sum of block gradient norms to the full gradient norm.
Since the blocks form a partition of the coordinates,
\[
\|\nabla f(\bm w^{t})\|^{2}
=\sum_{i=1}^{B}\|\nabla_{i}f(\bm w^{t})\|^{2}.
\tag{14}
\]
Insert (14) into (13):
\[
\E_{i_t}\!\big[f(\bm w^{t+1})\mid \bm w^{t}\big]
\le f(\bm w^{t})-\frac{\eta}{2B}\,\|\nabla f(\bm w^{t})\|^{2}.
\tag{15}
\]

\noindent\textbf{Step 4.}  Use Lemma~\ref{lem:gradbound} (8) to replace the gradient norm
by the suboptimality:
\[
\|\nabla f(\bm w^{t})\|^{2}
\ge 2L_{\max}\,\big(f(\bm w^{t})-f(\bm w^{\star})\big).
\tag{16}
\]
Plug (16) into (15) and recall $\eta=1/L_{\max}$:
\[
\E_{i_t}\!\big[f(\bm w^{t+1})\mid \bm w^{t}\big]
\le f(\bm w^{t})-\frac{1}{2B}\,
2\big(f(\bm w^{t})-f(\bm w^{\star})\big)
= f(\bm w^{t})-\frac{1}{B}\big(f(\bm w^{t})-f(\bm w^{\star})\big).
\tag{17}
\]

\noindent\textbf{Step 5.}  Rearrange (17) to isolate the expected suboptimality at
iteration $t+1$:
\[
\E_{i_t}\!\big[f(\bm w^{t+1})-f(\bm w^{\star})\mid \bm w^{t}\big]
\le \Big(1-\frac{1}{B}\Big)\big(f(\bm w^{t})-f(\bm w^{\star})\big).
\tag{18}
\]

\noindent\textbf{Step 6.}  Unroll the recursion.  Taking total expectation and
applying (18) repeatedly yields for any $t\ge 0$,
\[
\E\!\big[f(\bm w^{t})-f(\bm w^{\star})\big]
\le \Big(1-\frac{1}{B}\Big)^{t}\,
\big(f(\bm w^{0})-f(\bm w^{\star})\big).
\tag{19}
\]

\noindent\textbf{Step 7.}  The bound (19) is exponential, but we prefer a
$O(1/t)$ bound that does not rely on strong convexity.  From (12) we also have
the deterministic inequality
\[
f(\bm w^{t+1})-f(\bm w^{\star})
\le f(\bm w^{t})-f(\bm w^{\star})
-\frac{\eta}{2}\|\nabla_{i_t}f(\bm w^{t})\|^{2}.
\tag{20}
\]
Summing (20) from $t=0$ to $T-1$ and taking total expectation gives
\[
\E\!\big[f(\bm w^{T})-f(\bm w^{\star})\big]
\le f(\bm w^{0})-f(\bm w^{\star})
-\frac{\eta}{2}\sum_{t=0}^{T-1}
\E\!\big[\|\nabla_{i_t}f(\bm w^{t})\|^{2}\big].
\tag{21}
\]

\noindent\textbf{Step 8.}  Apply Lemma~\ref{lem:expgrad} to the expectation inside the
sum:
\[
\E\!\big[\|\nabla_{i_t}f(\bm w^{t})\|^{2}\big]
=\frac{1}{B}\,\E\!\big[\|\nabla f(\bm w^{t})\|^{2}\big].
\tag{22}
\]

\noindent\textbf{Step 9.}  Use Lemma~\ref{lem:gradbound} (8) in reverse,
$\|\nabla f(\bm w^{t})\|^{2}\ge 2L_{\max}\big(f(\bm w^{t})-f(\bm w^{\star})\big)$,
to obtain
\[
\E\!\big[\|\nabla_{i_t}f(\bm w^{t})\|^{2}\big]
\ge \frac{2L_{\max}}{B}\,
\E\!\big[f(\bm w^{t})-f(\bm w^{\star})\big].
\tag{23}
\]

\noindent\textbf{Step 10.}  Insert (23) into (21) and use $\eta=1/L_{\max}$:
\[
\E\!\big[f(\bm w^{T})-f(\bm w^{\star})\big]
\le f(\bm w^{0})-f(\bm w^{\star})
-\frac{1}{2L_{\max}}\sum_{t=0}^{T-1}
\frac{2L_{\max}}{B}\,
\E\!\big[f(\bm w^{t})-f(\bm w^{\star})\big].
\tag{24}
\]

\noindent\textbf{Step 11.}  Simplify (24):
\[
\E\!\big[f(\bm w^{T})-f(\bm w^{\star})\big]
\le f(\bm w^{0})-f(\bm w^{\star})
-\frac{1}{B}\sum_{t=0}^{T-1}
\E\!\big[f(\bm w^{t})-f(\bm w^{\star})\big].
\tag{25}
\]

\noindent\textbf{Step 12.}  Rearrange (25) to isolate the sum of expectations:
\[
\frac{1}{B}\sum_{t=0}^{T-1}
\E\!\big[f(\bm w^{t})-f(\bm w^{\star})\big]
\le f(\bm w^{0})-f(\bm w^{\star})
-\E\!\big[f(\bm w^{T})-f(\bm w^{\star})\big]
\le f(\bm w^{0})-f(\bm w^{\star}).
\tag{26}
\]

\noindent\textbf{Step 13.}  Since the sequence $\{f(\bm w^{t})\}$ is non‑increasing in
expectation, the largest term among $\{ \E[f(\bm w^{t})-f(\bm w^{\star})]\}_{t=0}^{T-1}$
is the first one.  Hence each term is bounded by the average:
\[
\E\!\big[f(\bm w^{T})-f(\bm w^{\star})\big]
\le \frac{B}{T}\,
\big(f(\bm w^{0})-f(\bm w^{\star})\big).
\tag{27}
\]

\noindent\textbf{Step 14.}  Relate the initial suboptimality to the distance
to the optimum using smoothness (1) with $\bm y=\bm w^{\star}$:
\[
f(\bm w^{0})-f(\bm w^{\star})
\le \frac{L_{\max}}{2}\,\|\bm w^{0}-\bm w^{\star}\|^{2}.
\tag{28}
\]

\noindent\textbf{Step 15.}  Substitute (28) into (27) and obtain the final bound:
\[
\E\!\big[f(\bm w^{T})-f(\bm w^{\star})\big]
\le \frac{B\,L_{\max}}{2T}\,
\|\bm w^{0}-\bm w^{\star}\|^{2},
\]
which is exactly the statement (4) of Theorem~\ref{thm:main}. ∎

\vspace{0.5em}
\section*{Rate Derivation (With Constants)}
The bound derived in Step 15 can be written as
\[
\boxed{
\E\!\big[f(\bm w^{T})-f(\bm w^{\star})\big]
\le
\underbrace{\frac{B\,L_{\max}}{2}}_{C}\,
\frac{\|\bm w^{0}-\bm w^{\star}\|^{2}}{T}
}
\]
where the constant $C$ depends linearly on the number of blocks $B$ and the
maximum block‑smoothness $L_{\max}$.  Thus the convergence rate is
$\mathcal{O}(1/T)$.

\vspace{0.5em}
\section*{Special Cases}
\begin{itemize}
\item \textbf{Single block ($B=1$).}  RBCD reduces to ordinary gradient descent
with stepsize $1/L_{\max}$ and the bound (4) becomes
$\frac{L_{\max}\|\bm w^{0}-\bm w^{\star}\|^{2}}{2T}$, the classical result.
\item \textbf{Uniform smoothness.}  If all blocks share the same Lipschitz constant
$L_i\equiv L$, then $L_{\max}=L$ and the bound simplifies to
$\frac{B L}{2T}\|\bm w^{0}-\bm w^{\star}\|^{2}$.
\item \textbf{Strong convexity.}  Adding the assumption
$f$ is $\mu$‑strongly convex yields a linear rate
\[
\E\!\big[f(\bm w^{T})-f(\bm w^{\star})\big]
\le
\big(1-\tfrac{\mu}{B L_{\max}}\big)^{T}
\big(f(\bm w^{0})-f(\bm w^{\star})\big).
\]
The proof follows the same steps with the additional inequality
$\frac{1}{2L_{\max}}\|\nabla f(\bm w)\|^{2}\ge \mu\big(f(\bm w)-f(\bm w^{\star})\big)$.
\end{itemize}

\end{document}